\begin{ex}
Đường gấp khúc ABCD gồm bao nhiêu đoạn thẳng?
\shortans{3}
\loigiai{
Gồm 3 đoạn thẳng: AB, BC, CD.
}
\end{ex}

\begin{ex}
Hình tứ giác có bao nhiêu đỉnh?
\shortans{4}
\loigiai{
Có 4 đỉnh.
}
\end{ex}

\begin{ex}
Hình tam giác có bao nhiêu cạnh?
\shortans{3}
\loigiai{
Có 3 cạnh.
}
\end{ex}

\begin{ex}
Khối trụ có bao nhiêu mặt phẳng đáy?
\shortans{2}
\loigiai{
Có 2 mặt đáy phẳng.
}
\end{ex}

\begin{ex}
Khối cầu có bao nhiêu đỉnh?
\shortans{0}
\loigiai{
Khối cầu có 0 đỉnh.
}
\end{ex}

\begin{ex}
Khối lập phương có bao nhiêu mặt vuông?
\shortans{6}
\loigiai{
Có 6 mặt vuông.
}
\end{ex}

\begin{ex}
Khối hộp chữ nhật có bao nhiêu đỉnh?
\shortans{8}
\loigiai{
Có 8 đỉnh.
}
\end{ex}

\begin{ex}
Khối hộp chữ nhật có bao nhiêu cạnh?
\shortans{12}
\loigiai{
Có 12 cạnh.
}
\end{ex}

\begin{ex}
Đường gấp khúc MNPQ gồm 3 đoạn thẳng dài lần lượt 10 cm, 20 cm, 30 cm. Độ dài đường gấp khúc là bao nhiêu xăng-ti-mét?
\shortans{60}
\loigiai{
Độ dài: $10 + 20 + 30 = 60$ cm.
}
\end{ex}

\begin{ex}
Chu vi hình tam giác có 3 cạnh dài lần lượt 12 cm, 15 cm, 18 cm là bao nhiêu xăng-ti-mét?
\shortans{45}
\loigiai{
Chu vi: $12 + 15 + 18 = 45$ cm.
}
\end{ex}

\begin{ex}
Chu vi hình tứ giác có độ dài các cạnh lần lượt là 10 cm, 15 cm, 20 cm, 25 cm là bao nhiêu xăng-ti-mét?
\shortans{70}
\loigiai{
Chu vi: $10 + 15 + 20 + 25 = 70$ cm.
}
\end{ex}

\begin{ex}
Một hình vuông có cạnh dài 9 cm. Chu vi hình vuông đó là bao nhiêu xăng-ti-mét?
\shortans{36}
\loigiai{
Chu vi: $9 \times 4 = 36$ cm.
}
\end{ex}

\begin{ex}
Một hình chữ nhật có chiều dài 15 cm, chiều rộng 10 cm. Chu vi hình chữ nhật đó là bao nhiêu xăng-ti-mét?
\shortans{50}
\loigiai{
Chu vi: $(15 + 10) \times 2 = 50$ cm.
}
\end{ex}

\begin{ex}
Cho 3 điểm A, B, C thẳng hàng theo thứ tự đó. Biết AB dài 14 cm, BC dài 16 cm. Độ dài đoạn AC là bao nhiêu xăng-ti-mét?
\shortans{30}
\loigiai{
Độ dài AC: $14 + 16 = 30$ cm.
}
\end{ex}

\begin{ex}
Đoạn thẳng MN dài 25 cm, điểm P nằm giữa M và N sao cho MP dài 10 cm. Đoạn PN dài bao nhiêu xăng-ti-mét?
\shortans{15}
\loigiai{
Đoạn PN dài: $25 - 10 = 15$ cm.
}
\end{ex}

\begin{ex}
Cần ít nhất bao nhiêu que tính để xếp được 1 hình vuông?
\shortans{4}
\loigiai{
Cần 4 que tính.
}
\end{ex}

\begin{ex}
Cần ít nhất bao nhiêu que tính để xếp được 1 hình tam giác?
\shortans{3}
\loigiai{
Cần 3 que tính.
}
\end{ex}

\begin{ex}
Dùng 15 que tính bằng nhau xếp được nhiều nhất bao nhiêu hình tam giác rời nhau?
\shortans{5}
\loigiai{
Xếp được: $15 : 3 = 5$ hình tam giác.
}
\end{ex}

\begin{ex}
Dùng 16 que tính bằng nhau xếp được nhiều nhất bao nhiêu hình vuông rời nhau?
\shortans{4}
\loigiai{
Xếp được: $16 : 4 = 4$ hình vuông.
}
\end{ex}

\begin{ex}
Hai hình tam giác rời nhau có tổng số cạnh là bao nhiêu?
\shortans{6}
\loigiai{
Có: $3 + 3 = 6$ cạnh.
}
\end{ex}

\begin{ex}
Hai hình vuông rời nhau có tổng số góc vuông là bao nhiêu?
\shortans{8}
\loigiai{
Có: $4 + 4 = 8$ góc vuông.
}
\end{ex}

\begin{ex}
Một hình tam giác có chu vi 36 cm, 3 cạnh có độ dài bằng nhau. Mỗi cạnh dài bao nhiêu xăng-ti-mét?
\shortans{12}
\loigiai{
Mỗi cạnh dài: $36 : 3 = 12$ cm.
}
\end{ex}

\begin{ex}
Một hình tứ giác có chu vi 60 cm, 4 cạnh có độ dài bằng nhau. Mỗi cạnh dài bao nhiêu xăng-ti-mét?
\shortans{15}
\loigiai{
Mỗi cạnh dài: $60 : 4 = 15$ cm.
}
\end{ex}

\begin{ex}
Tổng số mặt phẳng đáy của 2 khối trụ rời nhau là bao nhiêu?
\shortans{4}
\loigiai{
Mỗi khối trụ có 2 mặt đáy phẳng, hai khối có: $2 + 2 = 4$ mặt đáy phẳng.
}
\end{ex}

\begin{ex}
Hình chữ nhật có mấy cạnh?
\shortans{4}
\loigiai{
Hình chữ nhật có 4 cạnh.
}
\end{ex}
